\documentclass[11pt,a4paper]{article}
\usepackage[margin=2.2cm]{geometry}
\usepackage{times}
\usepackage{microtype}
\usepackage[hidelinks]{hyperref}
\usepackage{array}
\usepackage{booktabs}
\usepackage{tabularx}
\usepackage{enumitem}
\usepackage{titlesec}
\setlist{itemsep=1pt,topsep=3pt}
% Nama file dan kueri berhuruf mesin ketik tidak bisa dipenggal; biarkan spasi melebar daripada meluber.
\sloppy
\setlength{\parindent}{0pt}
\setlength{\parskip}{4pt}
\titleformat{\section}{\normalfont\Large\bfseries}{\thesection.}{0.5em}{}
\titlespacing*{\section}{0pt}{14pt}{6pt}
\titleformat{\subsection}{\normalfont\large\bfseries}{}{0em}{}
\titlespacing*{\subsection}{0pt}{10pt}{4pt}
\newcolumntype{Y}{>{\raggedright\arraybackslash}X}
\newcolumntype{P}[1]{>{\raggedright\arraybackslash}p{#1}}
\newcommand{\f}[1]{\texttt{#1}}
\newcommand{\tanya}[2]{\item \textit{``#1''}\\ #2}
% Dibuat otomatis oleh tools/latex_data.py dari output pipeline. Jangan diedit manual.
\newcommand{\ambang}{86,90}
\newcommand{\badDates}{0}
\newcommand{\bestMacro}{rf}
\newcommand{\bestModel}{Random Forest}
\newcommand{\boletoShare}{19,89}
\newcommand{\cardAvgPay}{165,81}
\newcommand{\cardInstallments}{3,51}
\newcommand{\cardShare}{77,00}
\newcommand{\cleanBadDate}{0}
\newcommand{\cleanGeoDup}{261.831}
\newcommand{\cleanGeoOut}{33}
\newcommand{\cleanPayment}{3}
\newcommand{\cleanReview}{551}
\newcommand{\cleanStatus}{1.234}
\newcommand{\cmRR}{8.729}
\newcommand{\cmRT}{1.186}
\newcommand{\cmTR}{751}
\newcommand{\cmTT}{9.176}
\newcommand{\curatedColumns}{35}
\newcommand{\cvFoneBest}{0,8339}
\newcommand{\cvFoneDepthTwo}{0,6970}
\newcommand{\cvFoneFifteen}{0,8339}
\newcommand{\cvFoneMinFifty}{0,8156}
\newcommand{\cvFoneTwentyFive}{0,8304}
\newcommand{\cvSeconds}{136,6}
\newcommand{\dtAcc}{0,8455}
\newcommand{\dtAuc}{0,8722}
\newcommand{\dtCvFone}{0,8339}
\newcommand{\dtDepth}{15}
\newcommand{\dtFit}{3,3}
\newcommand{\dtFmRendah}{0,8451}
\newcommand{\dtFmTinggi}{0,8460}
\newcommand{\dtFone}{0,8455}
\newcommand{\dtGap}{8,5}
\newcommand{\dtMinInst}{1}
\newcommand{\dtNodes}{6.333}
\newcommand{\dtPrecRendah}{0,8469}
\newcommand{\dtPrecTinggi}{0,8441}
\newcommand{\dtRecRendah}{0,8433}
\newcommand{\dtRecTinggi}{0,8478}
\newcommand{\dtTrainFone}{0,9304}
\newcommand{\dtbAcc}{0,8311}
\newcommand{\dtbAuc}{0,8679}
\newcommand{\dtbCvFone}{--}
\newcommand{\dtbFit}{3,1}
\newcommand{\dtbFmRendah}{0,8297}
\newcommand{\dtbFmTinggi}{0,8325}
\newcommand{\dtbFone}{0,8311}
\newcommand{\dtbPrecRendah}{0,8361}
\newcommand{\dtbPrecTinggi}{0,8262}
\newcommand{\dtbRecRendah}{0,8233}
\newcommand{\dtbRecTinggi}{0,8388}
\newcommand{\dtbTrainFone}{0,9210}
\newcommand{\dtdAcc}{0,7528}
\newcommand{\dtdAuc}{0,7924}
\newcommand{\dtdCvFone}{0,7535}
\newcommand{\dtdFit}{1,4}
\newcommand{\dtdFmRendah}{0,7614}
\newcommand{\dtdFmTinggi}{0,7436}
\newcommand{\dtdFone}{0,7525}
\newcommand{\dtdPrecRendah}{0,7354}
\newcommand{\dtdPrecTinggi}{0,7729}
\newcommand{\dtdRecRendah}{0,7893}
\newcommand{\dtdRecTinggi}{0,7163}
\newcommand{\dtdTrainFone}{0,7536}
\newcommand{\etlSeconds}{39,8}
\newcommand{\firstMonth}{2016-09}
\newcommand{\gainFeatures}{1,4}
\newcommand{\gainGbt}{4,3}
\newcommand{\gainRf}{5,7}
\newcommand{\gbtAcc}{0,8881}
\newcommand{\gbtAuc}{0,9570}
\newcommand{\gbtCvFone}{--}
\newcommand{\gbtFit}{61,9}
\newcommand{\gbtFmRendah}{0,8877}
\newcommand{\gbtFmTinggi}{0,8885}
\newcommand{\gbtFone}{0,8881}
\newcommand{\gbtPrecRendah}{0,8903}
\newcommand{\gbtPrecTinggi}{0,8860}
\newcommand{\gbtRecRendah}{0,8852}
\newcommand{\gbtRecTinggi}{0,8910}
\newcommand{\gbtTrainFone}{0,9139}
\newcommand{\geoAfterClean}{738.299}
\newcommand{\impA}{\texttt{category}}
\newcommand{\impB}{\texttt{total\_weight\_g}}
\newcommand{\impC}{\texttt{total\_freight}}
\newcommand{\impD}{\texttt{freight\_per\_item}}
\newcommand{\impE}{\texttt{avg\_weight\_g}}
\newcommand{\impValA}{0,1238}
\newcommand{\impValB}{0,1178}
\newcommand{\impValC}{0,1014}
\newcommand{\impValD}{0,0704}
\newcommand{\impValE}{0,0675}
\newcommand{\lastMonth}{2018-09}
\newcommand{\lateBestDays}{11,9}
\newcommand{\lateBestKm}{486}
\newcommand{\lateBestPct}{4,04}
\newcommand{\lateBestState}{PR}
\newcommand{\lateLow}{62,42}
\newcommand{\lateOrders}{6.381}
\newcommand{\latePct}{6,66}
\newcommand{\lateReview}{2,27}
\newcommand{\lateWorstDays}{21,5}
\newcommand{\lateWorstKm}{2.102}
\newcommand{\lateWorstPct}{17,43}
\newcommand{\lateWorstState}{MA}
\newcommand{\mostOrdersCategory}{\texttt{bed\_bath\_table}}
\newcommand{\nCustomers}{94.983}
\newcommand{\nLatih}{78.357}
\newcommand{\nOrders}{98.199}
\newcommand{\nRendah}{49.080}
\newcommand{\nStatesAnalysed}{17}
\newcommand{\nTinggi}{49.119}
\newcommand{\nUji}{19.842}
\newcommand{\onTimeLow}{9,27}
\newcommand{\onTimeOrders}{89.443}
\newcommand{\onTimeReview}{4,29}
\newcommand{\oneTimePct}{96,96}
\newcommand{\ordersAfterStatus}{98.207}
\newcommand{\ordersBeforeClean}{99.441}
\newcommand{\peakMonth}{2017-11}
\newcommand{\peakOrders}{7.421}
\newcommand{\peakRevenue}{1.003.862,14}
\newcommand{\rawRows}{1.550.922}
\newcommand{\rawSizeMB}{120,3}
\newcommand{\repeatPct}{3,04}
\newcommand{\rfAcc}{0,9024}
\newcommand{\rfAuc}{0,9674}
\newcommand{\rfCvFone}{--}
\newcommand{\rfFit}{103,2}
\newcommand{\rfFmRendah}{0,9001}
\newcommand{\rfFmTinggi}{0,9045}
\newcommand{\rfFone}{0,9023}
\newcommand{\rfPrecRendah}{0,9208}
\newcommand{\rfPrecTinggi}{0,8855}
\newcommand{\rfRecRendah}{0,8804}
\newcommand{\rfRecTinggi}{0,9243}
\newcommand{\rfTimesSlower}{31}
\newcommand{\rfTrainFone}{0,9978}
\newcommand{\rowsCategoryTranslation}{71}
\newcommand{\rowsCustomers}{99.441}
\newcommand{\rowsGeolocation}{1.000.163}
\newcommand{\rowsOrderItems}{112.650}
\newcommand{\rowsOrders}{99.441}
\newcommand{\rowsPayments}{103.886}
\newcommand{\rowsProducts}{32.951}
\newcommand{\rowsReviews}{99.224}
\newcommand{\rowsSellers}{3.095}
\newcommand{\scaleFoneFull}{0,8455}
\newcommand{\scaleFoneHalf}{0,8157}
\newcommand{\scaleFoneMin}{0,7595}
\newcommand{\scaleRatioFour}{2,4}
\newcommand{\scaleRowsMax}{313.428}
\newcommand{\scaleTimeDouble}{5,64}
\newcommand{\scaleTimeFull}{3,74}
\newcommand{\scaleTimeHalf}{3,04}
\newcommand{\scaleTimeMax}{9,07}
\newcommand{\scaleTimeMin}{2,01}
\newcommand{\secondCategory}{\texttt{watches\_gifts}}
\newcommand{\shallowNodes}{19}
\newcommand{\topCategory}{\texttt{health\_beauty}}
\newcommand{\topCategoryRevenue}{1.255.695,13}
\newcommand{\totalRevenue}{13.494.400,74}


\begin{document}

{\LARGE\bfseries Materi Presentasi Kelompok 1\par}
\vspace{2pt}
{\large Klasifikasi Nilai Transaksi E-Commerce Olist dengan HDFS, Spark, dan MLlib\par}
\vspace{4pt}
Anggota 1: Mohammad Ariffansyah (24051204099) \quad Anggota 2: [Nama] \quad
Anggota 3: [Nama] \quad Anggota 4: [Nama]

\section{Pembagian tugas dan urutan presentasi}
Urutan pembicara mengikuti nomor slide (\f{slides/presentasi.html}), jadi tidak ada
lompat-lompat. Daftar apa yang dijelaskan di setiap slide ada di Bagian~\ref{sec:naskah}.

\begin{tabularx}{\textwidth}{P{2.1cm}P{1.9cm}YP{4.6cm}}
\toprule
Siapa & Slide & Bagian & Demo langsung \\
\midrule
Anggota 1 & 1--6 & Pembuka, pertanyaan, alasan klasifikasi, dataset, arsitektur, HDFS & \f{./run.sh hdfs}, HDFS UI \f{localhost:9870} \\
Anggota 2 & 7--8 & Big data dan Spark ETL & \f{./run.sh etl}, Spark UI \f{localhost:4040} selama job berjalan \\
Anggota 3 & 9--10 & Analisis Spark SQL & --- \\
Anggota 4 & 11--14 & Machine learning: setup, hasil, aturan pohon, skala data & grafik di slide \\
Anggota 3 & 15 & MongoDB dan aplikasi demo & aplikasi \f{localhost:3000} \\
Anggota 1 & 16 & Kesimpulan, lalu tanya jawab & --- \\
\bottomrule
\end{tabularx}

\textbf{Persiapan sebelum presentasi} (tahap \f{ml} sekitar 8 menit, jangan dijalankan saat presentasi):
\begin{enumerate}
\item \f{./run.sh up}, lalu \f{./run.sh hdfs}, \f{./run.sh etl}, dan \f{./run.sh ml}.
\item Di terminal lain: \f{./run.sh app}, lalu buka \f{http://localhost:3000}.
\item Buka tab: HDFS UI (\f{localhost:9870} $\rightarrow$ Utilities $\rightarrow$ Browse $\rightarrow$
  \f{/olist}), aplikasi demo, folder \f{output/}, dan \f{artikel/artikel.pdf}.
\end{enumerate}

%=====================================================================
\section{Di mana bagian big data-nya}
Jawaban singkat: dari sisi \textbf{ukuran}, data ini belum termasuk big data. Yang
``big data'' adalah \textbf{cara data disimpan dan diolah}: terdistribusi, paralel, dan
bisa diperbesar tanpa mengubah kode.

\subsection{Teknologinya}
\begin{tabularx}{\textwidth}{P{2.3cm}YY}
\toprule
Alat & Peran big data & Wujudnya di proyek \\
\midrule
HDFS & Berkas dipecah menjadi blok dan disebar ke datanode. Kapasitas ditambah dengan
menambah node, bukan membeli disk yang lebih besar. &
9 CSV (\rawSizeMB{}~MB) di \f{/olist/raw}; hasil olahan Parquet di \f{/olist/curated};
model dan prediksi di \f{/olist/models} dan \f{/olist/predictions}. \\
Apache Spark & Data dibagi menjadi partisi yang diproses paralel dan di memori. Kode yang
sama jalan di laptop maupun di kluster ratusan node. &
Join 9 tabel, cleaning, dan agregasi dikerjakan 12 \emph{thread} sekaligus. \\
Spark MLlib & Model dilatih langsung di atas data terdistribusi, tanpa memuat semuanya
ke satu mesin. &
Cross-validation 27 kombinasi, Random Forest, dan GBT. \\
\bottomrule
\end{tabularx}

\textbf{Bukti angka untuk presentasi} (Spark UI, tab Executors, saat tahap MLlib): Spark
menjalankan \textbf{41.205 task}, membaca total \textbf{20,6~GiB} data, dan memindahkan
\textbf{3,8~GiB} data antar-partisi (\emph{shuffle}). Volumenya sebesar itu karena data
diproses berulang kali: 27 kombinasi $\times$ 3 fold, ditambah lima model.

\subsection{Karakteristik datanya (5V)}
\begin{tabularx}{\textwidth}{P{2.1cm}Y}
\toprule
V & Kondisi di proyek \\
\midrule
Volume & \rawRows{} baris mentah; tabel geolokasi saja \rowsGeolocation{} baris. Kecil
untuk ukuran industri, dan sudah ditulis jujur sebagai keterbatasan. \\
Variety & 9 tabel berbeda: transaksi, teks ulasan, koordinat GPS, tanggal, kategori produk. \\
Veracity & Data kotor yang sudah ditangani: \cleanGeoDup{} baris duplikat, koordinat di
luar Brasil, ulasan ganda, nilai kosong. \\
Value & Wawasan bisnis (pesanan terlambat mendapat ulasan \lateReview{} vs
\onTimeReview{}) dan model dengan F1 \rfFone. \\
Velocity & Tidak ada: prosesnya \emph{batch}, bukan \emph{streaming}. Bisa disebut sebagai
pengembangan berikutnya. \\
\bottomrule
\end{tabularx}

\textbf{Kalau dosen bertanya ``ini big data atau bukan?''}: datanya berukuran menengah,
tetapi seluruh alurnya dirancang supaya bisa diperbesar. Eksperimen skala sudah menguji
data latih sampai \scaleRowsMax{} baris, dan waktunya naik lebih lambat daripada jumlah
data. Untuk data 100 kali lebih besar cukup menambah datanode dan executor; kodenya tidak
berubah.

%=====================================================================
\section{Fungsi Spark SQL}
Spark SQL adalah \textbf{mesin kueri terdistribusi}. Kita menulis SQL biasa, lalu Spark
(lewat \emph{optimizer} Catalyst) mengubahnya menjadi job paralel di atas data di HDFS.
Bedanya dengan MySQL: data tidak harus muat di satu server, dan kueri dikerjakan banyak
core atau node sekaligus.
\begin{itemize}
\item \textbf{Analisis} (ketentuan minimal 3, kita membuat 6): tabel Parquet dibaca dari
  HDFS, didaftarkan sebagai \emph{view} \f{orders} dan \f{items}, lalu dikueri. Isinya
  tren bulanan, kategori teratas, keterlambatan per negara bagian, keterlambatan vs
  ulasan, metode pembayaran, dan jumlah pesanan per pelanggan.
\item \textbf{Fungsi SQL yang dipakai}:
  \begin{itemize}
  \item \f{GROUP BY} dengan \f{COUNT}, \f{SUM}, dan \f{AVG} untuk agregasi.
  \item \f{JOIN} untuk menggabungkan item dengan pesanan.
  \item \f{HAVING COUNT(*) >= 500} untuk membuang kelompok yang terlalu kecil.
  \item \emph{Window function} \f{SUM(COUNT(*)) OVER ()} untuk menghitung persentase.
  \item \f{date\_format} untuk mengelompokkan per bulan.
  \item \f{ORDER BY ... LIMIT} untuk mengambil 10 teratas.
  \end{itemize}
\item \textbf{Di tahap ETL} dipakai DataFrame API (\f{join}, \f{groupBy}, \f{filter},
  dan \emph{window} \f{row\_number} untuk memilih ulasan terakhir). Mesinnya sama dengan
  Spark SQL, hanya ditulis dengan Python.
\end{itemize}

%=====================================================================
\section{Fungsi MongoDB dan aplikasi demo}
HDFS cocok untuk menyimpan \textbf{berkas besar yang dibaca sekaligus} oleh Spark, tetapi
lambat kalau aplikasi hanya perlu satu data, misalnya ``prediksi untuk pesanan X''. MongoDB,
basis data NoSQL berbasis dokumen JSON, mengisi peran itu: \textbf{menyajikan hasil} supaya
bisa dibaca aplikasi dengan cepat tanpa menjalankan Spark.
\begin{itemize}
\item \f{olist.predictions}: \nUji{} dokumen hasil prediksi \bestModel{} untuk data uji,
  dengan format seperti contoh di panduan (\f{id\_transaksi}, \f{kategori\_prediksi},
  \f{kategori\_aktual}, \f{prob\_tinggi}, \f{model}, \f{sumber}).
\item \f{olist.metrics}: satu dokumen ringkasan metrik ditambahkan setiap kali model
  dilatih, sehingga riwayat semua percobaan tersimpan.
\item \textbf{Tulis}: Spark menulis langsung lewat MongoDB Spark Connector 10.4.1, yang
  versinya cocok dengan Spark 3.5 dan Scala 2.12.
\item \textbf{Baca}: Spark membaca kembali di akhir job; \f{./run.sh mongo} menjalankan
  \f{mongosh}; aplikasi demo membaca dengan driver MongoDB untuk Node.js.
\item Di panduan MongoDB adalah komponen \textbf{opsional (nilai tambah)}, bukan pengganti
  HDFS. HDFS untuk data mentah dan olahan; MongoDB untuk hasil akhir yang siap dipakai
  aplikasi.
\end{itemize}

\subsection{Aplikasi demo (\f{app-demo/}, Next.js)}
Jalankan dengan \f{./run.sh app}, buka \f{http://localhost:3000}. Setiap panel
menampilkan kueri yang dijalankan di sebelah hasilnya.

\begin{tabularx}{\textwidth}{P{3.2cm}P{4.4cm}Y}
\toprule
Bagian & Kueri & Yang ditunjukkan \\
\midrule
Cek transaksi & \f{findOne}, \f{\$sample} & Cari satu pesanan:
prediksi, kategori aktual, probabilitas, dan dokumen aslinya. Klik contoh ID acak. \\
Aktual vs prediksi & \f{aggregate(\$group)} & Confusion matrix dan akurasi
(\rfAcc), dihitung MongoDB dari semua dokumen. \\
Seberapa yakin model & \f{aggregate(\$bucket)} & Jumlah pesanan per rentang
probabilitas: sebagian besar dekat 0 atau 1, artinya model yakin. \\
Riwayat pelatihan & \f{find().\allowbreak sort().\allowbreak limit()} & Setiap kali \f{./run.sh ml}
dijalankan, satu dokumen metrik bertambah. \\
Hasil Spark SQL & SQL dari \f{etl\_analysis.py} & Keenam kueri Spark SQL dan tabel
hasilnya (\f{output/*.csv}). \\
\bottomrule
\end{tabularx}

\textbf{Alur demo yang disarankan}: klik satu contoh ID, lalu jelaskan bahwa data ini
ditulis Spark dan dibaca aplikasi tanpa Spark. Turun ke confusion matrix dan tunjukkan
bahwa akurasinya sama dengan di artikel. Terakhir, tunjukkan satu kueri Spark SQL beserta
hasilnya.

%=====================================================================
\section{Yang harus dipelajari setiap anggota}

\subsection{Anggota 1: pembuka, dataset, lingkungan, HDFS, kesimpulan}
\textbf{Baca:} artikel Bab 1, Subbab 2.1, 2.2, 3.7, Bab 4; file \f{compose.yaml},
\f{hadoop.env}, \f{run.sh}.
\begin{itemize}
\item Namenode vs datanode, blok, replikasi, \emph{safe mode}. \f{dfs.replication=1}
  karena hanya ada satu datanode.
\item Mode \f{local[*]}: Spark memakai 12 \emph{thread} dalam satu container, tetapi tetap
  membaca dan menulis lewat \f{hdfs://namenode:8020}.
\item Lisensi dataset CC BY-NC-SA 4.0 dan alasan memilih Olist: 9 tabel yang harus
  digabung.
\item Label Pembelian Tinggi/Rendah dan alasan memilih klasifikasi (\oneTimePct\% pelanggan
  hanya belanja sekali), karena ini dijelaskan di slide 2--3.
\end{itemize}
\textbf{Pertanyaan yang mungkin muncul:}
\begin{itemize}
\tanya{Kenapa pakai HDFS kalau cuma satu laptop?}{Arsitekturnya sama dengan kluster;
  tinggal menambah datanode. Spark membaca lewat alamat HDFS, bukan file lokal.}
\tanya{Kenapa pakai container?}{Supaya siapa pun bisa menjalankan ulang dengan satu
  perintah dan mendapat hasil yang sama.}
\end{itemize}

\subsection{Anggota 2: big data dan Spark ETL}
\textbf{Baca:} artikel Subbab 2.3 dan 3.2; Bagian 2 dokumen ini; \f{jobs/etl\_analysis.py} bagian 1--3.
\begin{itemize}
\item Opsi \f{inferSchema}, \f{multiLine}, \f{escape}: komentar ulasan berisi baris
  baru, jadi tanpa opsi ini satu ulasan pecah menjadi beberapa baris.
\item Tabel kualitas data: sel kosong, \cleanGeoDup{} baris geolokasi duplikat, dan 814
  \f{review\_id} ganda. Setiap aturan cleaning beserta alasannya.
\item Join 9 tabel menjadi \nOrders{} pesanan (\curatedColumns{} kolom), rumus Haversine,
  rasio ongkir, dan alasan menyimpan dalam format Parquet.
\item Dasar Spark: transformasi vs aksi, \emph{lazy evaluation}, \f{cache()}.
\end{itemize}
\textbf{Pertanyaan yang mungkin muncul:}
\begin{itemize}
\tanya{Kenapa nilai kosong diisi di MLlib, bukan di ETL?}{Supaya median dihitung dari data
  latih saja dan tidak bocor dari data uji.}
\tanya{Kenapa kategori produk diambil dari item pertama?}{Kalau diambil dari item termahal,
  harga ikut menentukan fitur, padahal harga adalah pembentuk label.}
\end{itemize}

\subsection{Anggota 3: Spark SQL dan MongoDB}
\textbf{Baca:} artikel Subbab 3.3; Bagian 3 dan 4 dokumen ini; \f{etl\_analysis.py}
bagian 4; \f{app-demo/app/page.js}.
\begin{itemize}
\item Keenam kueri Spark SQL dan fungsi SQL-nya (Bagian 3).
\item Temuan utama:
  \begin{itemize}
  \item Puncak penjualan \peakMonth{} (Black Friday).
  \item Pesanan terlambat mendapat ulasan \lateReview{} vs \onTimeReview.
  \item \cardShare\% pesanan dibayar dengan kartu kredit.
  \item \oneTimePct\% pelanggan hanya berbelanja sekali, yang menjadi alasan memilih
    klasifikasi, bukan clustering.
  \end{itemize}
\item MongoDB:
  \begin{itemize}
  \item Isi kedua koleksi.
  \item Mode \f{overwrite} (prediksi) vs \f{append} (metrik).
  \item Operasi \f{findOne}, \f{aggregate} dengan \f{\$group}/\f{\$bucket}/\f{\$sample},
    serta \f{find().sort()}.
  \end{itemize}
\item Menjalankan aplikasi demo dan menjelaskan setiap panelnya.
\end{itemize}
\textbf{Pertanyaan yang mungkin muncul:}
\begin{itemize}
\tanya{Kenapa perlu MongoDB kalau sudah ada HDFS?}{HDFS untuk berkas besar yang dibaca Spark
  sekaligus; MongoDB untuk mencari satu data dengan cepat dari aplikasi.}
\tanya{Kenapa ada filter minimal 500 pesanan?}{Supaya persentase keterlambatan tidak
  dihitung dari sedikit data.}
\tanya{Apa bedanya Spark SQL dengan MySQL?}{Spark SQL memecah kueri menjadi job paralel di
  atas data terdistribusi; MySQL berjalan di satu server.}
\end{itemize}

\subsection{Anggota 4: machine learning dan skala data (bagian terberat, mulai paling awal)}
\textbf{Baca:} artikel Subbab 2.4, 2.5, 3.4, 3.5, 3.6; \f{jobs/train\_models.py};
\f{output/ml\_decision\_tree\_shallow.txt}.
\begin{itemize}
\item Label: total harga $\geq$ median (\ambang{}~BRL), sehingga kelas seimbang. Harga dan
  kolom yang baru diketahui setelah pesanan selesai dikeluarkan dari fitur (kebocoran
  label).
\item Pipeline: \f{StringIndexer} (\f{keep}), \f{Imputer}, \f{VectorAssembler},
  \f{maxBins=128} (karena ada 74 kategori).
\item Pembagian data latih/uji berbasis hash (\nLatih{} / \nUji{}), supaya hasil bisa
  diulang persis.
\item Cross-validation 27 kombinasi menghasilkan kedalaman \dtDepth. F1 latih
  \dtTrainFone{} vs uji \dtFone: ada \emph{overfitting}, tetapi membatasi pohon tidak
  memperbaiki hasil pada data baru.
\item Perbandingan model:
  \begin{itemize}
  \item Random Forest: F1 \rfFone, AUC \rfAuc.
  \item GBT: F1 \gbtFone.
  \item Decision Tree: F1 \dtFone.
  \item Decision Tree dangkal: F1 \dtdFone.
  \end{itemize}
  Random Forest sekitar \rfTimesSlower{} kali lebih lama dilatih, dan aturannya tidak bisa
  dibaca.
\item Definisi akurasi, presisi, recall, F1, dan AUC. Confusion matrix: \cmTT{} dan
  \cmRR{} benar; \cmTR{} dan \cmRT{} salah.
\item Cara membaca aturan pohon dangkal: berat sekitar 1,2~kg, ongkir, dan jumlah cicilan.
\item Eksperimen skala: \scaleTimeMin{} detik (10\%) $\rightarrow$ \scaleTimeFull{} detik
  (100\%) $\rightarrow$ \scaleTimeMax{} detik (4 kali data). Naik lebih lambat dari jumlah
  data karena ada biaya tetap dan proses paralel.
\end{itemize}
\textbf{Pertanyaan yang mungkin muncul:}
\begin{itemize}
\tanya{Harga kan sudah diketahui saat bayar, lalu gunanya apa?}{Untuk menjelaskan ciri
  transaksi bernilai tinggi, misalnya untuk menentukan promo ongkir atau program cicilan,
  bukan untuk meramal.}
\tanya{Kenapa Random Forest lebih bagus?}{Rata-rata dari banyak pohon yang dilatih pada
  sampel berbeda mengurangi variansi satu pohon, dan probabilitasnya lebih halus, sehingga
  AUC naik paling banyak.}
\tanya{Kenapa Random Forest dan GBT tidak di-tuning?}{Satu kali latih Random Forest sekitar
  \rfFit{} detik, jadi cross-validation penuh menambah sekitar 25 menit. Perbandingannya
  jadi menguntungkan Decision Tree.}
\tanya{Kenapa F1 untuk data 2 dan 4 kali tidak dilaporkan?}{Data latih hanya diduplikasi,
  tidak ada informasi baru, jadi yang diukur hanya waktunya.}
\end{itemize}

%=====================================================================
\section{Naskah per slide}\label{sec:naskah}
Setiap slide berisi: siapa yang membawakan, apa yang \textbf{dijelaskan} (poin yang
diucapkan), dan apa yang \textbf{harus dipahami} supaya bisa menjawab pertanyaan. Target
waktu sekitar 45--60 detik per slide.

\newcommand{\slidenote}[4]{%
  \par\noindent\textbf{Slide #1 · #2} \hfill \textit{#3}\par
  \vspace{1pt}#4\par\vspace{6pt}}
\newcommand{\jelas}{\textbf{Dijelaskan:}}
\newcommand{\paham}{\textbf{Harus dipahami:}}

\slidenote{1}{Judul}{Anggota 1}{
\jelas{} perkenalan kelompok dan topik dalam satu kalimat: ``Kami mengklasifikasikan
pesanan Olist menjadi Pembelian Tinggi atau Rendah dengan HDFS, Spark, MLlib, dan MongoDB.''\\
\paham{} urutan rute data di label: Kaggle $\rightarrow$ HDFS $\rightarrow$ Spark
$\rightarrow$ MLlib $\rightarrow$ MongoDB.}

\slidenote{2}{Pertanyaan penelitian}{Anggota 1}{
\jelas{} label dibuat dari median total harga (\ambang{}~BRL), sehingga kelasnya seimbang
(\nTinggi{} vs \nRendah); model harus menebak \emph{tanpa} melihat harga.\\
\paham{} kenapa median (kelas seimbang, jadi akurasi tidak bisa ``curang'' dengan menebak
satu kelas) dan apa itu kebocoran label. Kalau ditanya ``harga kan sudah diketahui?'':
tujuannya menjelaskan ciri transaksi bernilai tinggi, bukan meramal.}

\slidenote{3}{Kenapa klasifikasi, bukan clustering}{Anggota 1}{
\jelas{} \oneTimePct\% pelanggan hanya belanja sekali, jadi pola pembelian per pelanggan
sangat tipis; topik Kelompok~1 memberi dua pilihan dan kami memilih klasifikasi transaksi.\\
\paham{} angka ini hasil analisis Spark SQL A6. Penelitian lain (Zeinali dkk.) tetap
membuat segmentasi RFM, tetapi pada data ini frekuensi hampir selalu 1.}

\slidenote{4}{Dataset}{Anggota 1}{
\jelas{} 9 tabel, \rawRows{} baris, \rawSizeMB{}~MB, Brasil 2016--2018; tabel terbesar
adalah geolokasi; lisensi CC BY-NC-SA 4.0.\\
\paham{} isi singkat setiap tabel (pesanan, item, pembayaran, ulasan, pelanggan, penjual,
produk, geolokasi, terjemahan kategori) dan kenapa banyak tabel membuat Spark perlu: semua
harus digabung per pesanan.}

\slidenote{5}{Arsitektur}{Anggota 1}{
\jelas{} alur dua baris di slide, empat container Podman, satu perintah \f{./run.sh},
dan hasil yang sama setiap dijalankan ulang.\\
\paham{} peran tiap kotak: HDFS untuk menyimpan, Spark untuk mengolah, MLlib untuk melatih
model, MongoDB untuk menyajikan hasil, Next.js untuk menampilkan. Parquet adalah format
kolom yang menyimpan tipe data.}

\slidenote{6}{HDFS}{Anggota 1 · demo HDFS UI}{
\jelas{} tiga perintah HDFS (\f{mkdir}, \f{put}, \f{du}) dan buktinya di HDFS UI; 9 berkas,
\rawSizeMB{}~MB, replikasi 1.\\
\paham{} namenode menyimpan metadata, datanode menyimpan blok; replikasi 1 karena hanya ada
satu datanode; \emph{safe mode} saat namenode baru menyala.}

\slidenote{7}{Di mana big data-nya}{Anggota 2}{
\jelas{} yang ``big data'' adalah cara mengolahnya: HDFS, Spark, MLlib; bukti dari Spark UI
(41.205 task, 20,6~GiB dibaca, 3,8~GiB \emph{shuffle}, 12 core).\\
\paham{} arti partisi, task, dan \emph{shuffle}; tabel 5V di Bagian~2; jawaban jujur untuk
``ini big data atau bukan?''.}

\slidenote{8}{Spark ETL}{Anggota 2 · demo \f{./run.sh etl}}{
\jelas{} masalah data yang dibuang atau dirapikan (duplikat geolokasi, pesanan batal, ulasan
ganda, koordinat di luar Brasil) sampai menjadi \nOrders{} pesanan, \curatedColumns{} kolom.\\
\paham{} opsi \f{multiLine}/\f{escape}, Haversine, rasio ongkir, \emph{lazy evaluation},
dan kenapa nilai kosong diisi di MLlib, bukan di ETL.}

\slidenote{9}{Keterlambatan dan ulasan}{Anggota 3}{
\jelas{} ulasan rata-rata \onTimeReview{} (tepat waktu) vs \lateReview{} (terlambat); ulasan
buruk naik dari \onTimeLow\% ke \lateLow\%; keterlambatan tertinggi di \lateWorstState.\\
\paham{} kueri \f{GROUP BY} + \f{AVG}, alasan filter \f{HAVING COUNT(*) >= 500}, dan
\emph{window function} untuk persentase.}

\slidenote{10}{Tren penjualan dan cara bayar}{Anggota 3}{
\jelas{} pendapatan naik sepanjang 2017 dengan puncak \peakMonth{} (Black Friday); \cardShare\%
pesanan memakai kartu kredit dengan rata-rata \cardInstallments{} kali cicilan.\\
\paham{} \f{date\_format} untuk per bulan; kenapa bulan dengan kurang dari 100 pesanan tidak
digambar (terlihat seperti penjualan anjlok padahal datanya belum lengkap).}

\slidenote{11}{Setup machine learning}{Anggota 4}{
\jelas{} label, 22 fitur numerik + 4 kategorikal, kolom yang sengaja dibuang, dan pipeline
StringIndexer $\rightarrow$ Imputer $\rightarrow$ VectorAssembler $\rightarrow$ model.\\
\paham{} fungsi tiap tahap pipeline, \f{maxBins=128}, pembagian data berbasis hash
(\nLatih{} latih / \nUji{} uji), dan 3-fold cross-validation 27 kombinasi.}

\slidenote{12}{Hasil model}{Anggota 4}{
\jelas{} Random Forest terbaik (F1 \rfFone, AUC \rfAuc), lalu GBT \gbtFone{} dan Decision
Tree \dtFone; fitur rasio menambah \gainFeatures{} poin; Random Forest \rfTimesSlower{} kali
lebih lama dilatih.\\
\paham{} definisi akurasi, presisi, recall, F1, AUC; kenapa Random Forest lebih baik;
\emph{overfitting} Decision Tree (latih \dtTrainFone{} vs uji \dtFone); kenapa RF/GBT tidak
di-tuning.}

\slidenote{13}{Aturan pohon dangkal}{Anggota 4}{
\jelas{} baca tiga aturan (berat $\pm$1,2~kg, ongkir, cicilan, kategori); F1 hanya \dtdFone{}
tetapi aturannya bisa dijelaskan ke tim bisnis.\\
\paham{} cara membaca pohon dari atas ke bawah; \emph{trade-off} akurasi vs kemudahan
dijelaskan; kenapa pohon kecil tidak \emph{overfitting} (latih \dtdTrainFone).}

\slidenote{14}{Eksperimen skala data}{Anggota 4}{
\jelas{} data 40 kali lebih banyak hanya butuh waktu sekitar 4,5 kali (\scaleTimeMin{}
$\rightarrow$ \scaleTimeMax{} detik); F1 terus naik dari \scaleFoneMin{} ke \scaleFoneFull.\\
\paham{} kenapa kenaikan waktu lebih lambat dari data (biaya tetap, paralel) dan kenapa skala
2 dan 4 kali hanya mengukur waktu (data diduplikasi).}

\slidenote{15}{MongoDB dan aplikasi}{Anggota 3 · demo \f{localhost:3000}}{
\jelas{} prediksi Random Forest disimpan di \f{olist.predictions} (\nUji{} dokumen) dan dibaca
aplikasi tanpa Spark. Demo: cari satu ID, lalu confusion matrix (akurasi \rfAcc), lalu satu
kueri Spark SQL.\\
\paham{} beda HDFS dan MongoDB; operasi \f{findOne}, \f{aggregate}
(\f{\$group}, \f{\$bucket}, \f{\$sample}); konektor 10.4.1; mode \f{overwrite} vs \f{append}.}

\slidenote{16}{Kesimpulan}{Anggota 1}{
\jelas{} tiga poin: pipeline bisa diulang dengan satu perintah; Random Forest F1 \rfFone{}
dan Decision Tree \dtFone; keterbatasan (satu laptop, satu datanode, RF dan GBT belum
di-tuning). Lalu buka sesi tanya jawab.\\
\paham{} semua angka di tabel ``Angka penting'' dan rencana pengembangan: kluster beberapa
node, tuning RF/GBT, prediksi keterlambatan pengiriman.}

%=====================================================================
\section{Wajib dikuasai semua anggota}
\begin{itemize}
\item Menjelaskan alur lengkap dalam satu menit: Kaggle $\rightarrow$ HDFS $\rightarrow$
  Spark ETL $\rightarrow$ Spark SQL $\rightarrow$ MLlib $\rightarrow$ MongoDB $\rightarrow$
  aplikasi.
\item Alasan memilih klasifikasi transaksi dan cara label dibuat.
\item Masalah yang ditemui dan solusinya (tabel di bawah); dosen suka menanyakan ini.
\item Mencoba \f{./run.sh} dan \f{./run.sh app} sendiri minimal sekali.
\end{itemize}

\begin{tabularx}{\textwidth}{P{5.2cm}Y}
\toprule
Masalah & Solusi \\
\midrule
Skrip menunggu selamanya saat cek datanode & \f{grep -q} memutus pipe lebih awal, podman
keluar dengan SIGPIPE, dan \f{pipefail} menganggapnya gagal. Diganti \f{grep >/dev/null}. \\
Hasil berubah setiap ETL dijalankan ulang & Pembagian data berbasis hash, lalu data
dipartisi dan diurutkan per \f{order\_id}. \\
Random Forest kehabisan memori & 12 \emph{thread} $\times$ 256~MB memori statistik
memenuhi \emph{heap}; \f{maxMemoryInMB} diturunkan ke 64. \\
Kedalaman terbaik berada di batas grid & Grid diperluas sampai 25, sehingga terlihat F1
sudah mendatar. \\
Kebocoran label & Harga, nilai pembayaran, ulasan, dan keterlambatan tidak dipakai sebagai
fitur. \\
\bottomrule
\end{tabularx}

\subsection{Angka penting}
\begin{tabularx}{\textwidth}{YY}
\toprule
Data & Model \\
\midrule
\rawRows{} baris mentah, 9 tabel, \rawSizeMB{}~MB & Random Forest: F1 \rfFone, AUC \rfAuc \\
\nOrders{} pesanan bersih, \curatedColumns{} kolom & Decision Tree: F1 \dtFone{} (kedalaman \dtDepth) \\
Ambang label \ambang{}~BRL (median) & GBT F1 \gbtFone; pohon dangkal F1 \dtdFone \\
\latePct\% pesanan terlambat; ulasan \lateReview{} vs \onTimeReview & Fitur rasio: +\gainFeatures{} poin F1 \\
\oneTimePct\% pelanggan hanya belanja sekali & Latih RF \rfFit{} dtk vs DT \dtFit{} dtk \\
\bottomrule
\end{tabularx}

\end{document}
